% Part: turing-machines
% Chapter: machines-computations
% Section: combining-machines

\documentclass[../../../include/open-logic-section]{subfiles}

\begin{document}

\olfileid[tr]{tur}{mac}{cmb}
\olsection{Turing Makinelerinin Birleştirilmesi}

\begin{explain}
Şimdiye dek gördüğümüz Turing makinesi örnekleri oldukça basitti. Oysa çağdaş
bir programlama diliyle çözülebilen her problem Turing makineleriyle de
çözülebilir. Daha karmaşık Turing makineleri kurabilmek için makineleri
birleştirebileceğimize kendimizi inandırmamız önemlidir; böylece bir süreci
daha basit parçalara ayırarak daha karmaşık problemleri çözen makineler
kurabiliriz. Karmaşık bir problemi bileşenlerine ayırmanın doğal bir yolunu
bulabilirsek problemi birkaç aşamada ele alabilir, birkaç basit Turing
makinesi oluşturup bunları problemi çözen tek bir makinede birleştirebiliriz.
Bu nokta, sonraki bölümde Durma Problemi ele alınırken özellikle önemlidir.

$M=\tuple{Q,\Sigma,q_0,\delta}$ ile
$M'=\tuple{Q',\Sigma',q_0',\delta'}$ Turing makinelerini nasıl birleştiririz?
$M$ durduktan sonraki şerit konfigürasyonunu, $M'$ makinesinin bir
çalışmasının girdi konfigürasyonu olarak kullanırız. Bunu yapan tek bir
$M\frown M'$ Turing makinesi elde etmek için şunları yapın:
\begin{enumerate}
    \item $M$ ile $M'$ ortak hiçbir duruma sahip olmayacak
    ($Q\cap Q'=\emptyset$) biçimde $M'$'nin bütün $Q'$ durumlarını yeniden
    numaralandırın (ya da yeniden etiketleyin).
    \item $M\frown M'$ makinesinin durumları $Q\cup Q'$ olsun.
    \item Şerit alfabesi $\Sigma\cup\Sigma'$ olsun.
    \item Başlangıç durumu $q_0$ olsun.
    \item Geçiş fonksiyonu aşağıdaki $\delta''$ fonksiyonu olsun:
    \[\delta''(q,\sigma) =
    \begin{cases}
      \delta(q,\sigma) & \text{$q \in Q$ ve $\delta(q,\sigma)$ tanımlıysa}\\
      \delta'(q,\sigma) & \text{$q \in Q'$ ise}\\
      \tuple{q_0', \sigma, \TMstay} & \text{$q \in Q$ ve
      $\delta(q,\sigma)$ tanımsızsa}
    \end{cases}\]
\end{enumerate}
Ortaya çıkan makine $q\in Q$ durumundayken $M$'nin, $q\in Q'$ durumundayken
$M'$'nin yönergelerini kullanır. $q\in Q$ durumundayken $M$'nin geçişinin
olmadığı bir $\sigma$ simgesini taradığında (yani $M$ duracakken) $M'$'nin
başlangıç durumuna geçer ve şeridin içeriğiyle kafa konumunu değiştirmez.

$M$ makinesi disiplinli değilse $M$ durduğunda şerit kafasının nerede
olduğunu bilmediğimize, dolayısıyla $M$'nin durma konfigürasyonunda kafanın
$1$. kareyi taramak zorunda olmadığına dikkat edin. Makineleri birleştirirken
bunu göz önünde bulundurmak önemlidir.
\end{explain}

\begin{ex}
\emph{Makinelerin Birleştirilmesi:} Uzunlukları $n$ ve~$m$ olan iki
$\TMstroke$ bloğundan oluşan girdide başlatıldığında şeritte $2(m+n)$ tane
$\TMstroke$ simgesinden oluşan tek bir blokla duran bir makine tasarlayacağız.
Bu makineyi kurmak için zaten bildiğimiz iki makineyi, toplama makinesiyle iki
katına çıkarıcıyı birleştirebiliriz. Önce toplama makinesinin durum diyagramını
çizelim.
\[
\begin{tikzpicture}[->,>=stealth',shorten >=1pt,auto,node distance=2.8cm,
                    semithick]
  \tikzstyle{every state}=[fill=none,draw=black,text=black]

  \node[initial,state] (A)              {$q_0$};
  \node[state]         (B) [right of=A] {$q_1$};
  \node[state]         (C) [right of=B] {$q_2$};

  \path (A) edge node {\TMtrans{\TMblank}{\TMstroke}{\TMstay}} (B)
            edge [loop above] node {\TMtrans{\TMstroke}{\TMstroke}{\TMright}} (B)
        (B) edge [loop above] node {\TMtrans{\TMstroke}{\TMstroke}{\TMright}} (B)
            edge node {\TMtrans{\TMblank}{\TMblank}{\TMleft}} (C)
        (C) edge [loop above] node {\TMtrans{\TMstroke}{\TMblank}{\TMstay}} (C);
\end{tikzpicture}
\]
$q_2$ durumunda durmak yerine çıktıyı iki katına çıkarmak için çalışmayı
sürdürmek istiyoruz. İki katına çıkarıcı makinenin girdideki ilk çizgiyi
sildiğini ve ayrı bir çıktıya iki çizgi yazdığını anımsayın. Şerit kafasının
toplama makinesi çıktısının ilk çizgisini okuduğundan emin olmak için bir
yönerge ekleyelim.
\[
\begin{tikzpicture}[->,>=stealth',shorten >=1pt,auto,node distance=2.8cm,
                    semithick]
  \tikzstyle{every state}=[fill=none,draw=black,text=black]

  \node[initial,state] (A)              {$q_0$};
  \node[state]         (B) [right of=A] {$q_1$};
  \node[state]         (C) [right of=B] {$q_2$};
  \node[state]         (D) [below left of=C] {$q_3$};
  \node[state]         (E) [below left of=D] {$q_4$};

  \path (A) edge node {\TMtrans{\TMblank}{\TMstroke}{\TMstay}} (B)
            edge [loop above] node {\TMtrans{\TMstroke}{\TMstroke}{\TMright}} (B)
        (B) edge [loop above] node {\TMtrans{\TMstroke}{\TMstroke}{\TMright}} (B)
            edge node {\TMtrans{\TMblank}{\TMblank}{\TMleft}} (C)
        (C) edge node {\TMtrans{\TMstroke}{\TMblank}{\TMleft}} (D)
        (D) edge [loop left] node {\TMtrans{\TMstroke}{\TMstroke}{\TMleft}} (D)
            edge node[left, xshift=-2mm] {\TMtrans{\TMendtape}{\TMendtape}{\TMright}} (E);
\end{tikzpicture}
\]
Artık girdiyi iki katına çıkarmak kolaydır: Yapmamız gereken tek şey iki
katına çıkarıcı makineyi $q_4$ durumuna bağlamaktır. Bunun için iki katına
çıkarıcı makinenin durumlarını $q_0$ yerine $q_4$'ten başlayacak biçimde
yeniden adlandırmalıyız; böylece iki başlangıç durumumuz olmaz. Son diyagram
\olref{fig:combined} içindeki gibi görünmelidir.
\begin{figure}
\[
\begin{tikzpicture}[->,>=stealth',shorten >=1pt,auto,node distance=2.8cm,
                    semithick]
  \tikzstyle{every state}=[fill=none,draw=black,text=black]
  \node[initial,state] (A)              {$q_0$};
  \node[state]         (B) [right of=A] {$q_1$};
  \node[state]         (C) [right of=B] {$q_2$};
  \node[state]         (D) [below left of=C] {$q_3$};
  \node[state]         (E) [below left of=D] {$q_4$};
  \node[state]         (2) [right of=E] {$q_5$};
  \node[state]         (3) [right of=2] {$q_6$};
  \node[state]         (4) [below of=3] {$q_7$};
  \node[state]         (5) [left of=4]  {$q_8$};
  \node[state]         (6) [left of=5]  {$q_9$};

  \path (A) edge node {\TMtrans{\TMblank}{\TMstroke}{\TMstay}} (B)
            edge [loop above] node {\TMtrans{\TMstroke}{\TMstroke}{\TMright}} (B)
        (B) edge [loop above] node {\TMtrans{\TMstroke}{\TMstroke}{\TMright}} (B)
            edge node {\TMtrans{\TMblank}{\TMblank}{\TMleft}} (C)
        (C) edge node {\TMtrans{\TMstroke}{\TMblank}{\TMleft}} (D)
        (D) edge [loop left] node {\TMtrans{\TMstroke}{\TMstroke}{\TMleft}} (D)
            edge node[left, xshift=-2mm] {\TMtrans{\TMendtape}{\TMendtape}{\TMright}} (E)
    (E) edge node {\TMtrans{\TMstroke}{\TMblank}{\TMright}} (2)
    (2) edge [loop above] node {\TMtrans{\TMstroke}{\TMstroke}{\TMright}} (2)
      edge node {\TMtrans{\TMblank}{\TMblank}{\TMright}} (3)
    (3) edge [loop above] node {\TMtrans{\TMstroke}{\TMstroke}{\TMright}} (3)
        edge node {\TMtrans{\TMblank}{\TMstroke}{\TMright}} (4)
    (4) edge [loop below] node {\TMtrans{\TMblank}{\TMstroke}{\TMleft}} (4)
        edge node {\TMtrans{\TMstroke}{\TMstroke}{\TMleft}} (5)
    (5) edge [loop below]  node {\TMtrans{\TMstroke}{\TMstroke}{\TMleft}} (5)
        edge              node {\TMtrans{\TMblank}{\TMblank}{\TMleft}} (6)
    (6) edge [loop below] node {\TMtrans{\TMstroke}{\TMstroke}{\TMleft}} (6)
        edge              node {\TMtrans{\TMblank}{\TMblank}{\TMright}} (E);
\end{tikzpicture}
\]
\caption{Toplama ve iki katına çıkarma makinelerinin birleştirilmesi}
\ollabel{fig:combined}
\end{figure}
\end{ex}

\begin{prop}
    $M$ ve $M'$ disiplinli olup sırasıyla $f\colon\Nat^k\to\Nat$ ve
    $f'\colon\Nat\to\Nat$ fonksiyonlarını hesaplıyorsa $M\frown M'$ de
    disiplinlidir ve $\comp{f}{f'}$ fonksiyonunu hesaplar.
\end{prop}

\begin{proof}
    $M$ disiplinli olduğundan $f(n_1,\dots,n_k)=m$ çıktısıyla durduğunda kafa
    $1$. kareyi tarar. Şimdi $M'$'nin başlangıç durumuna girersek makine yine
    $1$. kareyi tararken $f'(m)$ çıktısıyla durur.
    \olref[dis]{defn:disciplined} tanımındaki öteki koşullar da sağlanır.
\end{proof}

\begin{prob}
    \cref{tur:mac:dis:prob:disc-succ} içindeki $M$ makinesini alıp
    $M\frown M$'yi kurarak $f(x)=x+2$ hesaplayan disiplinli bir Turing
    makinesi verin.
\end{prob}
\end{document}
